\scanpage{111}
\mspage{102}
% archive sheet number: 111

suffit que les $\unsure{\Phi|\varphi}\in D$ composantes le soient. Pour que $\Phi_D$ soit ouverte, il f. et s. que les $\Phi\deleted{\unsure{\subset}}\in D$ le soient.
% margin note written slanted
\margin{\unsure{Ces deux énoncés|Ces deux derniers} \unsure{deux} faits sont des faits de topologie générale dans $\unsure{\mathscr{M}|\mathcal{M}}$, muni de $\unsure{\underline{\Delta}|\trianglelefteq}$.}

Dans cette description des préfigures et des figures, la relation $\overset{\circ}{\ll}$ n'est pas intervenue encore, ni (explicitement) la relation $\mathrel{|\circ|}$ --- à sa place, c'est la relation plus explicite $\underset{\mathrm{pos}}{<}$ (ou simplement $<$) dans $\mathscr{M}$ qui nous a servi. Mais maintenant que je passe aux supports, il n'en sera pas de même.

Il faut ensuite donner une détermination explicite du support d'une préfigure quelconque. Cela se fera en plusieurs étapes

\textcircled{\scriptsize $1^{\circ}$} Énoncé ``\unsure{additif|relatif}'' : si \unsure{la} $\Phi$ a pour composantes connexes les $\Phi_{\unsure{j|i}}$, on a
\[
\operatorname{supp}^{\circ}\Phi=\bigcup_i\operatorname{supp}^{\circ}(\Phi_i)
\]
(réunion disjointe).

$2^{\circ})$ Détermination de $\operatorname{supp}^{\circ}\Phi$ pour $\Phi$ préfigure connexe. Pour ceci, trois étapes :
